\documentclass{article}
\usepackage[utf8]{inputenc}
\usepackage[T1]{fontenc}
\usepackage[margin=2cm]{geometry}
\usepackage{hyperref}
\sloppy
\begin{document}
\section*{Historical source}
Encyclopaedia Britannica scan, source row 1430. Parse 3.5 extracted the text. Super 3.5 VL corrected it from an enlarged crop of each block. It can still contain recognition errors.
Page headers, footers, and page numbers are removed.

\subsection*{Untitled continuation (E0121)}
\textit{Continues an entry from an earlier page.}

over the continent, they had become a great and powerful people; and from their opposite dispositions and manners, became very frequently enemies to the Egyptians. To one Salatis our author ascribes the destruction of Thebes in Upper Egypt, so much celebrated by Homer for its magnificence. But this certainly cannot be the case; for Homer wrote long after the time of Joseph: and we find that even then the Egyptians had the shepherds in abhorrence, in all probability because they had been grievously oppressed by them. Mr Bruce counts three invasions of these people; the first, that of Salatis already mentioned, who overthrew the first dynasty of Egyptian kings from Memphis, and destroyed Thebes: the second was that of Sabacco or So; for according to him this was not the name of a single prince, but of a people, and signifies shepherds: and the third, after the building of Memphis, where 240,000 of them were besieged as above mentioned. But accounts of this kind are evidently insufficient in the highest degree; for how is it possible that the third invasion, antecedent to the building of Jerusalem, could be posterior to the second, if the latter happened only in the days of Hezekiah?

In these early ages, however, it would seem that the kingdom of Egypt had been very powerful and its dominion very widely extended; since we find it said, that the Badrarians revolted from Ozymandyas, another Egyptian king of very high antiquity, and of whose wealth the most marvellous accounts are given.

After an unknown interval of time from this monarchy existed Sefistris. He was the first great warrior whose conquests are recorded with any degree of difficulties. In what age of the world he lived, is uncertain. Some chronologers, among whom is Sir Isaac Newton, are of opinion, that he is the Sefac or Shishak, who took Jerusalem in the reign of Rehoboam the son of Solomon. Others, however, place him much earlier; and Mr Whiston will have him to be the Pharaoh who refused to part with the Israelites, and was at last drowned in the Red Sea. Mr Bryant endeavours to prove that no such person ever existed; but that in his history, as well as that of many ancient heroes, we have an abridgement of that of the Cushites or Babylonians, who spread themselves over great part of the known world, and every where brought the people in subjection to them. His reign is reckoned the most extraordinary part of the Egyptian history; and the following seems to be the least fabulous account that can be got of it. The father of Sefistris was told in a dream, by the god Vulcan, that his son, who was then newly born, or perhaps still unborn, should be lord of the whole earth. His father, upon the credit of this vision, got together all the males in the land of Egypt that were born on the same day with Sefistris; appointed nurses and proper persons to take care of them, and had them treated like his own child; being persuaded that they who had been the constant companions of his youth would prove the most faithful ministers and soldiers. As they grew up, they were inured to laborious exercises; and, in particular, were never permitted to take any food till they had performed a course of 180 furlongs, upwards of 22 of our miles. When the old king imagined they were sufficiently educated in the martial way he designed them to follow, they were sent by way

of trial of their abilities against the Arabians. In this expedition Sefotris proved successful, and in the end subdued that people who had never before been conquered. He was sent to the westward, and conquered the greatest part of Africa; nor could he be stopped in his career till he arrived at the Atlantic ocean. Whilst he was on this expedition, his father died; and then Sefotris resolved to fulfil the prediction of Vulcan, by actually conquering the whole world. As he knew that this must take up a long time, he prepared for his journey in the best manner possible. The kingdom he divided into 36 provinces, and endeavoured to secure the affections of the people by gifts both of money and land. He forgave all who had been guilty of offences, and discharged the debts of all his soldiers. He then constituted his brother Arais the supreme regent; but forbade him to use the diadem, and commanded him to offer no injury to the queen or her children, and to abstain from the royal concubines. His army consisted of 600,000 foot, 24,000 horse, and 27,000 chariots. Besides the land-forces, he had at sea two mighty fleets; one, according to Diodorus, of 400 sail. Of these fleets, one was designed to make conquests in the west, and the other in the east; and therefore the one was built on the Mediterranean and the other on the Red Sea. The first of these conquered Cyprus, the coast of Phoenicia, and several of the islands called Cyclades; the other conquered all the coasts of the Red Sea; but its progress was stopped by shoals and difficult places which the navigators could not pass, so that he seems not to have made many conquests by sea.

With the land-forces Sesostris marched against the Ethiopians and Troglodytes; whom he overcame, and obliged them to pay him a tribute of gold, ebony, and ivory. From thence he proceeded as far as the promontory of Dira, which lay near the straits of Babelmandel, where he set up a pillar with an inscription in faced characters. He then marched on to the country where cinnamon grows, or at least to some country where cinnamon at that time was brought, probably some place in India; and here he in like manner set up pillars, which were to be seen for many ages after. As to his farther conquests, it is agreed by almost all authors of antiquity, that he over-ran and pillaged the whole continent of Asia, and some part of Europe. He crofted the Ganges, and erected pillars on its banks; and from thence he is said to have marched eastward to the very extremity of the Asiatic continent. Returning from thence, he invaded the Scythians and Thracians; but all authors do not agree that he conquered them. Some even affirm, that he was overthrown by them with great slaughter, and obliged to abandon a great part of his booty and military forces. But whether he had good or bad success in these parts, it is a common opinion that he settled a colony in Colchis. Herodotus, however, who gives the most particular account of the conquests of this monarch, does not say whether the colony was designedly planted by Sesostris; or whether part of his army loitered behind the rest, and took up their residence in that region. From his own knowledge, he asserts, that the inhabitants of that country were undoubtedly of Egyptian descent. This was evident from the personal resemblance they bore to the Egyp-


\section*{Source record}
\begin{description}
\item[Scan ID] \texttt{britannica:309f765b1ecc83e0a7f7762e44b4a927a7555443:train:1430}
\item[Image file] \texttt{page-1430-4405fefdf7b6.jpg}
\item[Image SHA-256] \texttt{4405fefdf7b68a3c6472787e25922ed5125f57c753e650518ecdaaf630f59566}
\item[Dataset] NationalLibraryOfScotland/encyclopaedia\_britannica\_illustrated at revision \texttt{309f765b1ecc83e0a7f7762e44b4a927a7555443} (CC0-1.0)
\item[Extraction] model \texttt{nvidia/nemotron-3.5-parse}, raw reply \texttt{extract/raw/page-1430-4405fefdf7b6.json}
\item[Raw reply SHA-256] \texttt{8f69ecb267f77d8c6eef2c38be9b3ceea394702407f3164fd996055f82c1238d}
\item[Correction] model \texttt{nvidia/Nemotron-3.5-Super-VL}, block crops, raw replies \texttt{correct/raw/crop\_nothink/page-1430-4405fefdf7b6-b*.json}; super\_vl: 5
\item[Running header] EGY [ 368 ] E G Y
\end{description}
\end{document}
